\begin{figure}[htbp]
\centering
\begin{tikzpicture}[x=.55mm,y=-.55mm,>={Stealth[length=1.8mm]},line width=.45pt]
\definecolor{ch05figblue}{RGB}{78,130,191}
\definecolor{ch05figred}{RGB}{237,33,36}
\draw[->] (130.509,529.079)--(310.051,530.398);
\draw[->] (210.755,591.015)--(210.755,462.372);
\filldraw[fill=ch05figred,draw=ch05figred] (185.333,504.197) circle[radius=1.1mm];
\filldraw[fill=ch05figblue,draw=ch05figblue] (185.333,553.429) circle[radius=1.1mm];
\filldraw[fill=ch05figblue,draw=ch05figblue] (237.489,502.080) circle[radius=1.1mm];
\filldraw[fill=ch05figblue,draw=ch05figblue] (261.850,478.258) circle[radius=1.1mm];
\filldraw[fill=ch05figblue,draw=ch05figblue] (164.149,578.310) circle[radius=1.1mm];
\filldraw[fill=ch05figred,draw=ch05figred] (164.149,478.787) circle[radius=1.1mm];
\filldraw[fill=ch05figred,draw=ch05figred] (236.177,554.742) circle[radius=1.1mm];
\filldraw[fill=ch05figred,draw=ch05figred] (261.596,578.310) circle[radius=1.1mm];
\filldraw[fill=ch05figred,draw=ch05figred] (297.609,508.432) circle[radius=1.1mm];
\end{tikzpicture}
\origfig{3}{红色点与蓝色点之间可以用线性分类器分开。}
\label{fig:ch05-feature-map}
\end{figure}
